\section{Linear Programming: Interior-Point Methods}

In section 14.1, the authors outline the theoretical foundations of primal-dual path following interior point algorithms.
They introduce the duality measure
\[
    \mu := \frac{x^Ts}{n}
\] 
and derive an algorithm whose iterates are feasible solutions to the problem, reducing the duality measure.
In other sources, \textit{Introduction to Linear Optimization}, they relate this method with adding a barrier penalty to the objective.
Cover this in more detail here.

\subsection{Sets}
The book introduces a number of sets important to the analysis and design of IPMs.
We repeat their definitions here

Feasible set:
\begin{equation}
        \mathcal{F} = \{ (x, \lambda, s) | Ax = b, A^T \lambda + s = c, (x, s) \geq 0 \}
\end{equation}

Strictly feasible set:
\begin{equation}
    \mathcal{F}^\mathit{0} = \{ (x, \lambda, s) | Ax = b, A^T \lambda + s = c, (x, s) > 0 \}
\end{equation}

The set of points, $(x_\tau, \lambda_\tau, s_\tau)$, which satisfy the following equations,
\begin{align*}
    A^T\lambda + s &= c\\
    Ax &= b\\
    x_i s_i = \tau, \quad \forall 1, \dots, n
\end{align*}
make up the central path, 
\begin{equation}
    \mathcal{C} = \{ (x_\tau, \lambda_\tau, s_\tau | \tau > 0) \}
\end{equation}

Finally, two neighborhoods of the central path are introduced
\begin{equation}
    \mathcal{N}_2(\theta) =  \{ (x, \lambda, s) \in \mathcal{F}^\mathit{0} \quad | \quad \| XSe - \mu e \|_2 \leq \theta \mu \}
    \label{eq:ipm_set_n2}
\end{equation}
\begin{equation}
    \mathcal{N}_{-\infty}(\gamma) =  \{ (x, \lambda, s) \in \mathcal{F}^\mathit{0} \quad | \quad x_i s_i \geq \gamma \mu \> \forall i \}
\end{equation}

\subsection{Solutions}
The exercises chosen for this chapter are 14.1, 14.2, 14.5, 14.6
\subsubsection{Exercise 14.1}
The system (14.4a) is the following 
\[
    F(x,\lambda, s) = \begin{pmatrix}
        A^T\lambda + s - c\\
        Ax - b\\
        XSe
    \end{pmatrix} = 0
\]
and along with the bounds, $(x, s) \geq 0$ it constitutes the KKT conditions.

For our problem, we write out the above quantities explicitly
\begin{align*}
    A &= (1, 1)\\
    b &= 1\\
    c &= (1, 0)^T
\end{align*}
which means that 
\[
    F = \begin{pmatrix}
        \lambda + s_1 - 1\\
        \lambda + s_2\\
        x_1 + x_2 - 1\\
        x_1s_1\\
        x_2s_2
    \end{pmatrix}
\]
Inserting the two points listed in the exercise gives the following results
\begin{enumerate}
    \item[\textit{Actual Solution}: ]
        \[
            F = \begin{pmatrix}
                0 + 1 - 1\\
                0 + 0\\
                0 + 1 - 1\\
                0 \cdot 1\\
                1 \cdot 0
            \end{pmatrix} = \mathbf{0}
        \]
    \item[\textit{Spurious Solution}: ]
        \[
            F = \begin{pmatrix}
                1 + 0 - 1\\
                1 - 1\\
                1 + 0 - 1\\
                1 \cdot 0\\
                0 \cdot -1
            \end{pmatrix} = \mathbf{0}
        \]
\end{enumerate}

\subsubsection{Exercise 14.2}

That $\mathcal{N}_2(\theta_1) \subset \mathcal{N}_2(\theta_2)$ is equivalent to saying

\[
 (x, \lambda, s) \in \mathcal{N}_2(\theta_1) \implies (x, \lambda, s) \in \mathcal{N}_2(\theta_2)
\]
By using \eqref{eq:ipm_set_n2}, we can translate this two
\[
    \| XSe - \mu e \|_2 \leq \theta_1 \mu \implies \| XSe - \mu e \|_2 \leq \theta_2 \mu
\]
which, of course, is only true when $\theta_1 \leq \theta_2$.

We proceed similarly to show that
\[
 (x, \lambda, s) \in \mathcal{N}_{-\infty}(\theta_1) \implies (x, \lambda, s) \in \mathcal{N}_{-\infty}(\theta_2).
\]
This is equivalent to
\[
   x_i s_i \geq \gamma_1 \mu \implies x_i s_i \geq \gamma_2 \mu
\]
which can only be true if $\gamma_1 \geq \gamma_2$

\subsubsection{Exercise 14.5}
To show that $\mathcal{N}_2(0) = \mathcal{N}_{-\infty}(1) = \mathcal{C}$, we show that the conditions 
\begin{equation}
    x_i s_i \geq \mu, \quad \forall i = 1,\dots,n
    \label{eq:145cond_one}
\end{equation}
and 
\begin{equation}
    \| XSe - \mu e \|_2 = 0
    \label{eq:145cond_two}    
\end{equation}
imply $x_i s_i = \mu, \quad \forall i = 1,\dots,n$.

Beginning with \eqref{eq:145cond_one}, we have that
\[
    \mu \leq x_i s_i = n \mu - \sum_{j \neq i} x_i s_i \leq \mu 
\]
And, since $i$ is arbitrary, $x_i s_i = \mu, \> \forall i$.
This means $\mathcal{N}_{-\infty}(1)$ is the set of feasible points where all pairwise products $x_i s_i$ are equal which is the central path.

Moving on to condition \eqref{eq:145cond_two}, we begin by noting that $XSe$ is the vector of the pairwise products $x_i s_i$.
The the $L_2$ norm, $d: x, y \mapsto \|x - y\|_2$ is a valid metric on $\mathbb{R}^n$ which means $d(x,y) = 0 \implies x = y$. 
Thus $d(XSe, \mu e) = 0 \implies x_i s_i = \mu$, which was to be shown.    
\subsubsection{Exercise 14.6}